
% ============================================================================
\clearpage
\section*{v4.8.0-r03 Source-Only: Dimensional and Numerical Audit of the Physical-Constant Layer}
\addcontentsline{toc}{section}{v4.8.0-r03 Source-Only: Dimensional and Numerical Audit of the Physical-Constant Layer}
\label{sec:v480-r03-dimensional-audit}
% ============================================================================

\paragraph{Purpose.}
An external structural evaluation of v4.7.0 rated equation and dimensional integrity as the
weakest of its five criteria.  This layer audits the places where the monolith assigns
physical units, CODATA values or frame-insertion counts: the nomenclature, the
phase-shifted constants of \S\ref{sec:phase-constants}, the Hawking analysis of
\S\ref{sec:bh}, the cosmological statements around
Theorems~\ref{thm:dark-offset} and~\ref{thm:dark-energy}, and the photon-frame
\(\Lambda\) block of \S\ref{sec:lambda_cancel}.  Every number was recomputed.  Earlier layers
are not rewritten; each finding is recorded here with the correction that follows from the
document's own rules, or with the reason why no such correction exists.

\noindent\textbf{Source-only status.}
This v4.8.0-r03 layer is source-only.  It does not build a PDF, does not assign a new DOI,
and does not authorize a public release.  The v4.8.0 Version DOI
\href{https://doi.org/10.5281/zenodo.23365244}{10.5281/zenodo.23365244} and the Concept DOI
\href{https://doi.org/10.5281/zenodo.19210728}{10.5281/zenodo.19210728} are unchanged by it.

\subsection*{v4.8.0-r03 the frame-insertion calculus}
\label{subsec:v480-r03-k-calculus}

The definition in \S\ref{sec:phase-constants} assigns to a constant \(X\) the photon-frame
value \(X^{(\mathrm{photon})}=\gL^{-k(X)}X^{(\mathrm{matter})}\).

\paragraph{Lemma v4.8.0-r03.1.}
If \(X=\prod_i Y_i^{a_i}\) with real exponents \(a_i\) and every factor obeys the rule, then
\(X\) obeys it with
\[
k(X)=\sum_i a_i\,k(Y_i).
\]

\paragraph{Proof.}
\(X^{(\mathrm{photon})}=\prod_i\bigl(\gL^{-k(Y_i)}Y_i^{(\mathrm{matter})}\bigr)^{a_i}
=\gL^{-\sum_i a_ik(Y_i)}X^{(\mathrm{matter})}\).  \(\square\)

The lemma is what Theorem~\ref{thm:planck-photon} uses: \(\lP=(\hbar G_N/c^3)^{1/2}\) gives
\(k(\lP)=\tfrac12(1+2-0)=\tfrac32\), as Proposition~\ref{prop:uv-photon} also states.  Since
the rule multiplies dimensionful quantities by the pure number \(\gL\), it is dimensionally
consistent by construction.  What it requires is that every composite quantity be assigned
the \(k\) the lemma gives, and that every factor of the composite carry its own \(k\).  The
findings below are the places where that is not done.

\subsection*{v4.8.0-r03 findings}
\label{subsec:v480-r03-findings}

\paragraph{F1 --- vacuum-energy ratio (\(k\)-calculus).}
Equation~\eqref{eq:rho-ratio}, Corollary~\ref{cor:matter-error} and the corollary
``Cosmological constant revisited'' in \S\ref{sec:phase-constants} use
\(\rho_{\mathrm{vac}}\sim\lP^{-4}\) and obtain the frame ratio \(\gL^{6}\).  The relation
\(\rho\sim\lP^{-4}\) holds only in units with \(\hbar=c=1\).  The framework assigns
\(k(\hbar)=1\), so \(\hbar\) cannot be set to one in both frames.  As an energy density,
\(\rho_{\mathrm{vac}}\sim\hbar c\,\lP^{-4}\), and Lemma~v4.8.0-r03.1 gives
\[
k(\rho_{\mathrm{vac}})=1+0-4\cdot\tfrac32=-5,
\qquad
\frac{\rho_{\mathrm{vac}}^{(\mathrm{photon})}}{\rho_{\mathrm{vac}}^{(\mathrm{matter})}}=\gL^{5}=22.53\ldots,
\]
not \(\gL^{6}\).  The mass-density form \(m_P/\lP^{3}\) gives the same \(k=-\tfrac12-\tfrac92=-5\).
In addition, the source quotes \(\gL^{6}\) as \(41.8\) in one place and \(42.0\) in another;
with \(\gL=1.86444\ldots\) the value is \(42.004\ldots\).

\paragraph{F2 --- fine-structure drift (sign of the exponent and empirical status).}
With \(k(\alpha^{-1})=1\) (Theorem~\ref{thm:k-values}) the rule gives
\((\alpha^{-1})^{(\mathrm{photon})}=\gL^{-1}(\alpha^{-1})^{(\mathrm{matter})}\), hence
\[
\alpha^{(\mathrm{matter})}=\gL^{-1}\,\alpha^{(\mathrm{photon})},
\]
whereas the corollary ``\(\alpha\)-drift as a frame effect'' writes
\(\alpha^{(\mathrm{matter})}(t)=\gL^{(\mathrm{local})}(t)\,\alpha^{(\mathrm{photon})}\).  The two
are inverse to each other.  Separately, the corollary speaks of an ``observed NIST drift''.  The
NIST Al\(^+\)/Hg\(^+\) clock comparison of 2008 found a drift consistent with zero, and a fit of
the same data for an annual modulation gives a fractional amplitude of about
\((-8.9\pm8.4)\times10^{-18}\), also consistent with zero.  There is therefore no observed drift
to explain.  The annual-modulation prediction survives only as a constraint:
\(|\Delta\gL^{(\mathrm{local})}/\gL|\lesssim10^{-17}\) over the year.

\paragraph{F3 --- three incompatible types for \(\Lambda\).}
The monolith uses \(\Lambda\) in three dimensionally different roles:
\[
\begin{array}{ll}
\Lambda_{\mathrm{bare}}\sim\gamma_1^2\approx200\,m_P^4 & (\text{Prop.~\ref{prop:bare-lambda}}),\\
\Lambda_{\mathrm{eff}}\sim\bsym^2\rho_P & (\text{near Thm.~\ref{thm:dark-offset}}),\\
\Lambda_{\mathrm{obs}}=\bsym^2\,(\lP^{(\mathrm{photon})})^{-2} & (\text{\S\ref{sec:falsifiable}}).
\end{array}
\]
The first equates a pure number with a Planck-unit energy density, the second is an energy
density, the third a curvature with unit \(\mathrm{m}^{-2}\).  Curvature and energy density are
related by \(\Lambda=8\pi G_N\rho_\Lambda/c^4\); the factor \(8\pi\) is absent throughout.  The
photon-frame step of Proposition~\ref{prop:bare-lambda} uses \(k(\Lambda)=2\), while the
curvature form \(\Lambda\propto\lP^{-2}\) gives \(k=-3\) under Lemma~v4.8.0-r03.1.  The source
already states that a self-consistent \(k(\Lambda)\) is open; r03 records that the value \(2\)
used in the display is not the one the calculus gives for any of the three forms.

\paragraph{F4 --- the ``unit-system artifact'' remark.}
Remark~\ref{rem:open} and the prediction P4 of \S\ref{sec:falsifiable} state that the
\(10^{-122}\) of the cosmological-constant problem is an artefact of SI units, and that in
natural units \(\Lambda_{\mathrm{obs}}\sim\bsym^2\approx0.30\).  This is a dimensional error.
The quantity
\[
\Lambda_{\mathrm{obs}}\,\lP^{2}=1.1056\times10^{-52}\,\mathrm{m}^{-2}\cdot(1.6163\times10^{-35}\,\mathrm{m})^{2}
=2.89\times10^{-122}
\]
is dimensionless and takes the same value in every unit system; in units with \(\lP=1\) the
measured \(\Lambda_{\mathrm{obs}}\) is \(2.89\times10^{-122}\), not \(0.30\).  Taken literally,
\(\Lambda_{\mathrm{obs}}=\bsym^2(\lP^{(\mathrm{photon})})^{-2}=7.5\times10^{69}\,\mathrm{m}^{-2}\)
exceeds the measured value by a factor \(6.8\times10^{121}\).  P4 is therefore contradicted by
measurement as stated.  The same block also gives three different sizes for the required
cancellation (\(2.3\times10^{-4}\), \(1.5\times10^{-3}\) and a factor \(3.1\)), which cannot all
describe one ratio.

\paragraph{F5 --- cosmological fractions.}
Theorem~\ref{thm:dark-energy} writes \(1-\bsym^2\approx\tfrac23\approx0.698\); with
\(\bsym^2=0.30178\ldots\) the value is \(0.69822\ldots\), which is not \(\tfrac23\).  Two different
predictions for \(\Omega_\Lambda/\Omega_M\) occur: \(\bsym^2/\gstring=2.37\) near
Theorem~\ref{thm:dark-offset} and \((1-\bsym^2)/\bsym^2=2.31\) from
Theorem~\ref{thm:dark-energy}.  They are different functions of the same parameters and cannot
both hold.  The observed ratio from the Planck 2018 values quoted in
Remark~\ref{rem:omega-bsym} is \(0.6889/0.3111=2.214\), not \(2.18\).  The same remark writes
\(\bsym=\gBI/\gstring\approx0.5493\); that quotient is \(\gL=1.864\), and the stated value
belongs to \(\bsym=\sqrt{(\gL-1)/(\gL+1)}\).

\paragraph{F6 --- Hawking temperature.}
Proposition~\ref{prop:hawking-gBI-tachyon} derives
\(T_H=\hbar c^3/(8\pi G_NMk_B)=m_P^2c^2/(8\pi Mk_B)\) correctly.  The step to
\(k(T_H)=-1\) declares \(k(k_B)=0\) but also needs \(k(M)=0\) for the black-hole mass, which is
not declared, and which differs from the assignment \(k(m_ec^2)=1\) for the electron.  Under
Lemma~v4.8.0-r03.1, \(k(T_H)=1-2-k(M)=-1-k(M)\).

\paragraph{F7 --- notation.}
In Theorem~\ref{thm:dark-offset} the letter \(k\) denotes a spring constant, while \(k(X)\)
denotes the frame-insertion count elsewhere in the same part.  The oscillator frequency
\(\omega=\gstring\sqrt{2/3}\) is a pure number; Planck units are implied but not declared.

\paragraph{Verified without finding.}
\(\gstring=\ln2/(\pi\sqrt3)=0.12738\ldots\);
\(\gL=\gBI/\gstring=1.86444\ldots\);
\(\bsym=\sqrt{(\gL-1)/(\gL+1)}=0.54935\ldots\);
\(\gL^{3/2}=2.5458\ldots\);
\(\lP^{(\mathrm{photon})}=6.35\times10^{-36}\,\mathrm{m}\);
\(k(\lP)=\tfrac32\);
the units of \(G_N\), \(\hbar\), \(c\) in Theorem~\ref{thm:k-values};
\(T_H=\hbar c^3/(8\pi G_NMk_B)\);
\((\alpha^{-1})^{(\mathrm{photon})}=73.50\ldots\) under \(k(\alpha^{-1})=1\).

\paragraph{Theoremlet v4.8.0-r03.}
Under the frame-insertion rule of \S\ref{sec:phase-constants} and Lemma~v4.8.0-r03.1:
\begin{enumerate}[leftmargin=2em]
 \item the photon-to-matter ratio of a Planck-cutoff vacuum energy density is \(\gL^{5}\), not
 \(\gL^{6}\);
 \item \(k(\alpha^{-1})=1\) implies \(\alpha^{(\mathrm{matter})}=\gL^{-1}\alpha^{(\mathrm{photon})}\);
 \item \(\Lambda\) has no single \(k\) until one of its three forms is chosen, and none of them
 gives \(k=2\);
 \item no choice of units turns the dimensionless number \(\Lambda_{\mathrm{obs}}\lP^2\approx
 2.9\times10^{-122}\) into a number of order one.
\end{enumerate}

\paragraph{Proof sketch.}
Items 1--3 are applications of the lemma to \(\hbar c\,\lP^{-4}\), to \(\alpha=(\alpha^{-1})^{-1}\),
and to \(\lP^{-2}\), \(\hbar c\,\lP^{-4}\) and \(G_N\rho/c^4\) respectively.  Item 4 holds because a
ratio of two quantities of the same dimension is invariant under change of units.  \(\square\)

\subsection*{v4.8.0-r03 re-check of P4}
\label{subsec:v480-r03-p4-recheck}

F4 was checked a second time, against the main paper and against a reading of P4 that is
dimensionally consistent.

\paragraph{Origin of the unit \(\lP^{-2}\).}
The main paper states the frame contribution without a unit.  Its Theorem~11.6 (Dark energy as
permanent frame offset) gives \(\Lambda_{\mathrm{eff}}\sim\bsym^2\rho_P\), its Proposition~12.4
(Frame term gives the cosmological constant) gives \(\Lambda_{\mathrm{obs}}\sim\bsym^2\), and its
remark on the cosmological-constant problem records that reducing the bare value by
\(10^{122}\) through the \(\Theta\)--\(R\) cancellation ``is not yet established''
\cite{brendecke2026main}.  The unit \(\lP^{-2}\) and the statement that \(10^{-122}\) is a unit
artefact enter only in the photon-frame \(\Lambda\) block of the companion,
\S\ref{sec:lambda_cancel}.  The error recorded in F4 therefore lies in that later
reinterpretation, not in the main paper.

\paragraph{F8 --- ``dimensionally correct''.}
The proof of Theorem~\ref{thm:lambda-zero}(iii) calls \(\Lambda_{\mathrm{obs}}=\bsym^2/\lP^{(\mathrm{matter})2}\)
``dimensionally correct up to the \(\gL^{3/2}\) Planck-length correction''.  The unit
\(\mathrm{m}^{-2}\) is correct; the value is \(1.2\times10^{69}\,\mathrm{m}^{-2}\), a factor
\(10^{121}\) above the measurement, and no power of \(\gL\) closes that gap.  The theorem uses the
matter-frame Planck length, P4 the photon-frame one.

\paragraph{F9 --- citation.}
The proof of Theorem~\ref{thm:lambda-zero}(iii) and the status table of the \(\Lambda\) block cite
``Proposition~18.4 of the main paper''.  The main paper has no numbered environment 18.4.  The
statement meant is its Proposition~12.4, which asserts only the proportionality
\(\Lambda_{\mathrm{obs}}\sim\bsym^2\).

\paragraph{F10 --- units of the balance condition.}
In Equation~\eqref{eq:tr-balance} the integrand \(\Sigma(\tau)\) contains \(e^{-\tau/2}\), so
\(\tau\) is a pure number, while the upper limit \(\lP^{(\mathrm{photon})}\) and the right-hand
side \(\bsym^2\lP^{(\mathrm{photon})2}\) are written as a length and a squared length.  The
condition is consistent only when read in matter-frame Planck units, where
\(\lP^{(\mathrm{photon})}=\gL^{-3/2}=0.3928\ldots\) and the right-hand side is
\(\bsym^2\gL^{-3}=0.0466\ldots\).  This reading is not stated in the source.

\paragraph{A dimensionally consistent reading, offered and not adopted.}
Since \(\bsym^2\) is a pure number, the consistent place for it is a dimensionless fraction of
the critical density, not \(\Lambda\) itself.  If the Markov statement of
Theorem~\ref{thm:dark-energy}, \(\Omega_\Lambda=1-\bsym^2\), is taken as the prediction, then
\[
\Lambda_{\mathrm{obs}}=\frac{3\,\Omega_\Lambda H_0^2}{c^2}
=\frac{3\,(1-\bsym^2)\,H_0^2}{c^2}
=1.112\times10^{-52}\,\mathrm{m}^{-2}
\qquad(H_0=67.4\ \mathrm{km\,s^{-1}\,Mpc^{-1}}),
\]
which differs from \(\Lambda_{\mathrm{obs}}=1.1056\times10^{-52}\,\mathrm{m}^{-2}\) by \(0.6\%\).
At the level of the fraction, \(1-\bsym^2=0.6982\ldots\) compares with the Planck 2018 values
\(0.6847\) and \(0.6889\) by \(2.0\%\) and \(1.4\%\).

The reading has a price.  It predicts the fraction \(\Omega_\Lambda\), and takes the absolute scale
from the measured \(H_0\).  The number \(10^{-122}\) does not disappear: it becomes
\(\Lambda_{\mathrm{obs}}\lP^2=3\Omega_\Lambda(\lP/R_H)^2\) with the Hubble radius \(R_H=c/H_0\), so
that the remaining question is the hierarchy \(R_H/\lP=8.49\times10^{60}\).  The source
quantity closest to it is the active zero count \(n_H\sim10^{62}\) at the Hubble time
(Theorem~11.5 of the main paper); it differs from \(R_H/\lP\) by a factor of about \(12\), and the
source does not identify the two.  The reading also selects \(1-\bsym^2\) among the two
conflicting predictions of F5.

Whether P4 is restated in this form, restated otherwise, or withdrawn is left to the author.
r03 records the reading as an option and asserts nothing beyond the arithmetic above.

\subsection*{v4.8.0-r03 provenance of the frame parameters}
\label{subsec:v480-r03-frame-parameter-provenance}

\paragraph{F11 --- both inputs of \(\gL\) are values of the same LQG parameter.}
The frame factor \(\gL=\gBI/\gstring=1.864\ldots\) enters every frame conversion of the monolith,
\(\bsym=\sqrt{(\gL-1)/(\gL+1)}\), the cosmological fractions of
Theorem~\ref{thm:dark-energy}, the photon-frame Planck length of
Theorem~\ref{thm:planck-photon} and the instance \(\lambda_\gamma=\gL^{-3/2}\) of
v4.8.0-r01.  Its two inputs have the following origin in the literature.
\begin{enumerate}[leftmargin=2em]
 \item \(\gstring=\ln2/(\pi\sqrt3)=0.12738\ldots\), introduced in the nomenclature as
 ``string coupling'', is the value of the Barbero--Immirzi parameter obtained by Ashtekar, Baez,
 Corichi and Krasnov from the state counting of quantum isolated horizons
 (arXiv \texttt{gr-qc/9710007}, \texttt{gr-qc/0005126}).  It is not a string-theory coupling.
 \item \(\gBI\approx0.2375\) is the value of the same parameter after the corrected state
 counting of Domagala and Lewandowski and of Meissner (arXiv \texttt{gr-qc/0407051},
 \texttt{gr-qc/0407052}).  Domagala and Lewandowski show that the earlier counting contains a
 spurious constraint that suppresses spins above \(j=\tfrac12\), and that the earlier value
 therefore does not hold.
\end{enumerate}
Hence \(\gL\) is the ratio of the corrected to the superseded determination of one parameter.
Theorem~\ref{thm:frame-offset} reads this ratio as a Lorentz factor between the matter and the
photon frame.  The source supplies no derivation of either input from the HC trace formula: the
proof of Theorem~\ref{thm:frame-offset} says that \(\gstring\) is computed in the photon frame
from that formula, while the proof of the OP~1/T4 statement around
Remark~\ref{rem:goe-gue-string-lqg} calls it extrinsically given.  The two statements are
incompatible, and only the second is supported by the text.

\paragraph{F11a --- the entropy factor.}
The remark ``The entropy factor'' in \S\ref{sec:bh} writes the LQG entropy as
\(S=A\ln2/(8\pi\gBI\lP^2)\), finds agreement with \(S_{BH}=A/(4\lP^2)\) at
\(\gBI=\ln2/(2\pi)\approx0.110\), and records a discrepancy factor \(2.15\) as an open quantitative
problem.  The entropy formula of the counting that produced \(\ln2/(\pi\sqrt3)\) is
\[
S=\frac{\ln2}{4\pi\sqrt3\,\gamma}\,\frac{A}{\lP^2},
\]
which agrees with \(S_{BH}\) exactly at \(\gamma=\ln2/(\pi\sqrt3)=\gstring\).  The factor \(2.15\)
is produced by the formula used in the remark and is not a property of the framework; the
discrepancy between the two countings is \(\gBI/\gstring=\gL\) itself.

\paragraph{Consequence.}
The agreement of \(1-\bsym^2\) with the measured \(\Omega_\Lambda\) at the \(1\)--\(2\%\) level
recorded in \S\ref{subsec:v480-r03-p4-recheck} rests on \(\bsym\), and \(\bsym\) rests on the ratio
of two countings of one parameter, one of which is superseded.  As long as that is not stated,
the agreement reads as better supported than it is.

\paragraph{A reading in the readout vocabulary, offered and not adopted.}
The superseded counting admits only minimal spins \(j=\tfrac12\) at the horizon punctures, the
corrected one all spins.  \(\gL\) could therefore be read as the ratio between a minimal-spin
readout and a full-spin readout of the same horizon, two readouts in the sense of r81--r82 and
r93.  Under that reading \(\gL\) is not a Lorentz factor, and the chain
\(\gL\to\bsym\to\Omega_\Lambda\) would need a new justification: why the cosmological fraction
should be governed by the ratio of two horizon readouts.  r03 records the reading as an option
and asserts nothing beyond the provenance above.  Whether \(\gstring\) is renamed, the chain is
rederived, or this reading is adopted is left to the author.

\subsection*{v4.8.0-r03 status of the affected statements}
\label{subsec:v480-r03-status}

\[
\begin{array}{l|l|l}
\text{statement} & \text{finding} & \text{status after r03}\\
\hline
\text{Eq.~\eqref{eq:rho-ratio}, Cor.~\ref{cor:matter-error}} & \text{F1} & \text{corrected: } \gL^5\approx22.5\\
\alpha\text{-drift corollary} & \text{F2} & \text{exponent corrected; no observed drift}\\
\Lambda\text{ in Prop.~\ref{prop:bare-lambda}, Thm.~\ref{thm:dark-offset}, P4} & \text{F3} & \text{open: one type for }\Lambda\text{ to be chosen}\\
\text{Rem.~\ref{rem:open}, P4} & \text{F4} & \text{contradicted as stated; }\Omega\text{-reading offered}\\
\text{Thm.~\ref{thm:lambda-zero}(iii)} & \text{F8, F9} & \text{value off by }10^{121}\text{; citation corrected}\\
\text{Eq.~\eqref{eq:tr-balance}} & \text{F10} & \text{Planck units made explicit}\\
\gstring,\ \text{Thm.~\ref{thm:frame-offset}} & \text{F11} & \text{provenance recorded; readout reading offered}\\
\text{entropy-factor remark, \S\ref{sec:bh}} & \text{F11a} & \text{corrected: no factor }2.15\\
\Omega\text{ statements, Rem.~\ref{rem:omega-bsym}} & \text{F5} & \text{values corrected; two predictions conflict}\\
\text{Prop.~\ref{prop:hawking-gBI-tachyon}} & \text{F6} & \text{assumption } k(M)=0 \text{ made explicit}\\
\text{Thm.~\ref{thm:dark-offset}} & \text{F7} & \text{notation only}
\end{array}
\]

None of these statements is one of the guarded claims of the release scope, and none of the
five guards depends on them.  F3 and F4 concern the photon-frame \(\Lambda\) programme, which the
source already marks as open; r03 sharpens that status.  For P4 the source-only layer records
the contradiction and does not supply a replacement; a reformulation of the
\(\Lambda\) prediction, or its withdrawal, is left to the author.

\subsection*{v4.8.0-r03 ledger and guard update}
\label{subsec:v480-r03-ledger-guard}

\[
R_{\rm dim}^{v480\text{-}r03}=(0,0,1,1,1,1)
\]
for (\(k\)-calculus corrections F1, F2, F6 derived; numerical corrections F1, F5 recorded; single
type for \(\Lambda\) not chosen; a replacement for P4 offered in \S\ref{subsec:v480-r03-p4-recheck} but not adopted; one of the two
\(\Omega_\Lambda/\Omega_M\) predictions not selected; independent justification of \(\gL\) as a frame factor not supplied, \S\ref{subsec:v480-r03-frame-parameter-provenance}).

The layer spends no guard:
\[
R_{\rm glob}^{(v480\text{-}r03,+)}=(0,1,1,1,1),
\qquad
\Gamma_{\rm guard}^{v480\text{-}r03}=(1,1,1,1,1),
\qquad
\beta_4=1/256\ \text{unspent}.
\]
It does not imply exact finite \(C5\) zero closure, exact all-readout terminality,
\(C6/T2\), a global photon-clock/RH theorem, or global all-readout closure.
